% Generated draft: requires mathematical and source review.
\documentclass[11pt]{article}
\usepackage[margin=25mm]{geometry}
\usepackage{fontspec}
\setmainfont{Latin Modern Roman}
\usepackage{amsmath,amssymb,mathtools,unicode-math}
\setmathfont{Latin Modern Math}
\usepackage{enumitem,needspace,xurl,hyperref}
\setlength{\parindent}{0pt}
\setlength{\parskip}{5pt}
\setlength{\emergencystretch}{3em}
\newcommand{\N}{\mathbb{N}}
\newcommand{\Z}{\mathbb{Z}}
\newcommand{\Q}{\mathbb{Q}}
\newcommand{\R}{\mathbb{R}}
\newcommand{\C}{\mathbb{C}}
\newcommand{\F}{\mathbb{F}}
\newcommand{\A}{\mathbb{A}}
\newcommand{\PP}{\mathbb{P}}
\newcommand{\E}{\mathbb{E}}
\newcommand{\eps}{\varepsilon}
\newcommand{\dd}{\,\mathrm{d}}
\begin{document}
\section*{UnsolvedMath ID: 30004262}
\textbf{Source problem code:} catalogue entry
\subsection{Definitions and mathematical statement}
Let \((X,d)\) be a compact metric space and let \(N\) be a fixed finite Euclidean dimension. Interpreting preservation of equality and inequality of distances as preservation of equality and disequality, a \emph{loose embedding} into \(\R^N\) is a continuous map \(f:X\to\R^N\) such that, for all \(x,y,u,v\in X\),
\[
 d(x,y)=d(u,v)
 \quad\Longleftrightarrow\quad
 \|f(x)-f(y)\|_2=\|f(u)-f(v)\|_2.
\]
This condition implies injectivity by comparison with a diagonal pair. The supplied wording does not explicitly specify whether strict ordering of distances must also be preserved; the displayed definition uses disequality, without imposing an order-preservation condition.

The finite-subspace hypothesis is
\[
 \text{for every finite }F\subseteq X,\quad
 \text{there exists a loose embedding }
 f_F:(F,d|_{F\times F})\longrightarrow\R^N.
\]
The dimension \(N\) is the same for every \(F\), but the maps \(f_F\) may depend on \(F\); no compatibility between these maps is assumed.

The question in the supplied clean statement is whether this hypothesis implies the existence of a loose embedding \(f:(X,d)\to\R^N\), for every compact metric space \((X,d)\) and every fixed finite dimension \(N\).

The accompanying manifold question, read according to the original statement's phrase ``satisfy this hypothesis,'' is whether every compact Riemannian manifold \((M,g)\), equipped with its geodesic distance \(d_g\), admits a finite dimension \(N=N(M,g)\) such that every finite subspace of \((M,d_g)\) loosely embeds into \(\R^N\). Here geodesic distance means the infimum of the Riemannian lengths of curves joining the two points. The record does not specify connectedness or whether manifolds with boundary are included. In particular, the usual geodesic distance between different connected components is not finite, so the intended class must be clarified before this is a fully specified metric-space assertion. The clean statement's phrase ``have this property'' could instead refer to loose embeddability of the whole manifold; that alternative target remains unresolved.

There is also a substantive discrepancy between the supplied formulations: the original statement asks whether \(X\) is \emph{uniformly embeddable} into \(\R^N\), whereas the clean statement asks whether it is \emph{loosely embeddable} there. The record supplies no definition of uniform embeddability or explanation identifying these targets. The displayed question follows the clean statement, which the registry labels \texttt{corrected\_verified}; the original target is retained as an unresolved alternative, and that registry label is not an independent verification.
\subsection{Short English statement}
If every finite subset of a compact metric space admits a continuous embedding into one fixed Euclidean dimension that preserves exactly which distances are equal, must the entire space admit such an embedding? Does every compact Riemannian manifold satisfy the finite-subset hypothesis for some fixed dimension? The source leaves unresolved the distinction between uniform and loose embeddability, the intended manifold conventions, and the referent of ``this property.''
\subsection{Sources}
\begin{itemize}
\item Local UnsolvedMath registry: supplied complete source record, permanent ID \(30004262\), exact source problem code \texttt{catalogue entry}, title ``Finite-Dimensional Uniform Embeddings of Compact Metric Spaces.'' The record attributes the problem to Oberwolfach Reports workshop discussions; no individual proposer is supplied. It records proposed year \(2019\), status \texttt{open}, and statement status \texttt{corrected\_verified}. These are supplied registry labels, not independently established facts.
\item Supplied report citation: \emph{Mini-Workshop: Operator Algebraic Quantum Groups (2020)}, Oberwolfach Reports, pp.\ 2821--2867, DOI \texttt{10.4171/owr/2019/45}: \url{https://doi.org/10.4171/owr/2019/45}. The citation's parenthetical year \(2020\) differs from the supplied proposed year \(2019\) and the year appearing in the DOI; the record does not explain this difference.
\item Supplied literature source, labeled ``OWR 2019 45'': \url{https://ems.press/content/serial-article-files/46823}.
\item Literature assessment dated exactly \texttt{2026-08-22}: the supplied curation labels the problem open, reports that loose embeddability of a constructed metric space and the corresponding compact Riemannian manifold question were left open, and states that no general resolution was found. This is supplied dated curation, not independently verified. The constructed metric space is not described in the supplied record.
\item The supplied statement-verification note says that the clean statement was checked against the official Oberwolfach report and preserves its hypotheses, notation, and open target. This claim has not been independently checked here; the formulation discrepancies above still require human review.
\end{itemize}
\end{document}
